% GENERATED FILE. Edit canonical lessons, then run npm run generate:books.
% canonical-source-hash: fnv1a64:fa89239456b82085
% canonical-lessons: ES-C703-atardecer, ES-C703-anochecer, ES-C703-horizonte, ES-C703-arco-iris, ES-C703-contemplar, ES-R703-repaso-696-703, ES-R703-repaso-final-703, ES-R703-repaso-problemas-de-viaje

\chapter{Al caer la tarde}
\label{ch:es-al-caer-la-tarde}
\hlchaptermodality{703}

\begin{chapteropening}
\textbf{By the end of this chapter:} \emph{I can describe the end of a day out: sunset, nightfall, the horizon and a rainbow.}

Complete the last lesson of chapter 703: I can describe the end of a day out: sunset, nightfall, the horizon and a rainbow.
\end{chapteropening}

\section[el atardecer]{el atardecer — dusk, sunset}

\label{lesson:ES-C703-atardecer}



\begin{quote}
Before the new one: say the Spanish for the mud, then the Spanish for the raincoat.
\end{quote}

\subsection*{You'll want to know: el atardecer}
\textbf{el atardecer} — \textquotedblleft{}dusk, sunset\textquotedblright{}.

\begin{quote}
El atardecer en la playa es precioso.
\end{quote}

\begin{quote}
\emph{Sunset on the beach is beautiful.}
\end{quote}

Also a verb: \emph{atardece} — evening is falling.

\begin{cousinweb}
\emph{atardecer} is built on \emph{tarde} \textquotedblleft{}afternoon, late\textquotedblright{}, which is inherited from Latin \emph{tardē} \textquotedblleft{}late\textquotedblright{} (EtymDB), from \emph{tardus} \textquotedblleft{}slow\textquotedblright{}.
\end{cousinweb}

\begin{grammarlens}[title={What you've built}]
One more word for telling, explaining and describing at B1.
\end{grammarlens}

\subsection*{Your turn}
\begin{itemize}
\raggedright
  \item \emph{Say it:} \emph{el atardecer}
  \item \emph{Say it:} \emph{el atardecer}, once more
  \item \emph{Say it:} a sentence of your own with \emph{el atardecer}
  \item \emph{Recall:} say the Spanish for the mud, then the Spanish for the raincoat, then say \emph{el atardecer} again
\end{itemize}

\begin{tcolorbox}[breakable,colback=teal!4,colframe=teal!35!black,title={Before you move on}]
What does el atardecer mean? (\textquotedblleft{}Dusk, sunset\textquotedblright{}.) Say it once more.
\end{tcolorbox}
\section[anochecer]{anochecer — to get dark; nightfall}

\label{lesson:ES-C703-anochecer}



\begin{quote}
Before the new one: say the Spanish for the raincoat, then the Spanish for dusk.
\end{quote}

\subsection*{You'll want to know: anochecer}
\textbf{anochecer} — \textquotedblleft{}to get dark; nightfall\textquotedblright{}.

\begin{quote}
En invierno anochece muy pronto.
\end{quote}

\begin{quote}
\emph{In winter it gets dark very early.}
\end{quote}

\emph{Al anochecer} — at nightfall.

\begin{cousinweb}
\emph{anochecer} is built from \emph{a-} and Latin \emph{noctēscere} \textquotedblleft{}to grow dark\textquotedblright{} (EtymDB), from \emph{nox, noctis} \textquotedblleft{}night\textquotedblright{}.
\end{cousinweb}

\begin{grammarlens}[title={What you've built}]
One more word for telling, explaining and describing at B1.
\end{grammarlens}

\subsection*{Your turn}
\begin{itemize}
\raggedright
  \item \emph{Say it:} \emph{anochecer}
  \item \emph{Say it:} \emph{anochecer}, once more
  \item \emph{Say it:} a sentence of your own with \emph{anochecer}
  \item \emph{Recall:} say the Spanish for the raincoat, then the Spanish for dusk, then say \emph{anochecer} again
\end{itemize}

\begin{tcolorbox}[breakable,colback=teal!4,colframe=teal!35!black,title={Before you move on}]
What does anochecer mean? (\textquotedblleft{}To get dark; nightfall\textquotedblright{}.) Say it once more.
\end{tcolorbox}
\section[el horizonte]{el horizonte — the horizon}

\label{lesson:ES-C703-horizonte}



\begin{quote}
Before the new one: say the Spanish for dusk, then the Spanish for to get dark.
\end{quote}

\subsection*{You'll want to know: el horizonte}
\textbf{el horizonte} — \textquotedblleft{}the horizon\textquotedblright{}.

\begin{quote}
El sol se pone en el horizonte.
\end{quote}

\begin{quote}
\emph{The sun sets on the horizon.}
\end{quote}

\emph{Nuevos horizontes} — new horizons.

\begin{cousinweb}
\emph{horizonte} is a learned borrowing from Latin \emph{horizōn}, from Greek \emph{horízōn} \textquotedblleft{}bounding (circle)\textquotedblright{} (EtymDB), from \emph{hóros} \textquotedblleft{}boundary\textquotedblright{}.
\end{cousinweb}

\begin{grammarlens}[title={What you've built}]
One more word for telling, explaining and describing at B1.
\end{grammarlens}

\subsection*{Your turn}
\begin{itemize}
\raggedright
  \item \emph{Say it:} \emph{el horizonte}
  \item \emph{Say it:} \emph{el horizonte}, once more
  \item \emph{Say it:} a sentence of your own with \emph{el horizonte}
  \item \emph{Recall:} say the Spanish for dusk, then the Spanish for to get dark, then say \emph{el horizonte} again
\end{itemize}

\begin{tcolorbox}[breakable,colback=teal!4,colframe=teal!35!black,title={Before you move on}]
What does el horizonte mean? (\textquotedblleft{}The horizon\textquotedblright{}.) Say it once more.
\end{tcolorbox}
\section[el arco iris]{el arco iris — the rainbow}

\label{lesson:ES-C703-arco-iris}



\begin{quote}
Before the new one: say the Spanish for to get dark, then the Spanish for the horizon.
\end{quote}

\subsection*{You'll want to know: el arco iris}
\textbf{el arco iris} — \textquotedblleft{}the rainbow\textquotedblright{}.

\begin{quote}
Después de la lluvia salió un arco iris.
\end{quote}

\begin{quote}
\emph{After the rain a rainbow came out.}
\end{quote}

Also written \emph{arcoíris}; the plural is \emph{los arcoíris}.

\begin{cousinweb}
\emph{arco iris} is built from \emph{arco} \textquotedblleft{}bow\textquotedblright{}, inherited from Latin \emph{arcus} (EtymDB), and \emph{iris} \textquotedblleft{}rainbow\textquotedblright{}, from Latin \emph{īris} and Greek \emph{îris} (EtymDB).
\end{cousinweb}

\begin{grammarlens}[title={What you've built}]
One more word for telling, explaining and describing at B1.
\end{grammarlens}

\subsection*{Your turn}
\begin{itemize}
\raggedright
  \item \emph{Say it:} \emph{el arco iris}
  \item \emph{Say it:} \emph{el arco iris}, once more
  \item \emph{Say it:} a sentence of your own with \emph{el arco iris}
  \item \emph{Recall:} say the Spanish for to get dark, then the Spanish for the horizon, then say \emph{el arco iris} again
\end{itemize}

\begin{tcolorbox}[breakable,colback=teal!4,colframe=teal!35!black,title={Before you move on}]
What does el arco iris mean? (\textquotedblleft{}The rainbow\textquotedblright{}.) Say it once more.
\end{tcolorbox}
\section[contemplar]{contemplar — to gaze at, to contemplate}

\label{lesson:ES-C703-contemplar}



\begin{quote}
Before the new one: say the Spanish for the horizon, then the Spanish for the rainbow.
\end{quote}

\subsection*{You'll want to know: contemplar}
\textbf{contemplar} — \textquotedblleft{}to gaze at, to contemplate\textquotedblright{}.

\begin{quote}
Contemplamos el paisaje en silencio.
\end{quote}

\begin{quote}
\emph{We gazed at the landscape in silence.}
\end{quote}

Slower and calmer than \emph{mirar}.

\begin{cousinweb}
\emph{contemplar} is a learned borrowing from Latin \emph{contemplārī} \textquotedblleft{}to observe, to gaze at\textquotedblright{} (EtymDB), built on \emph{templum}, an open space marked out for watching the sky.
\end{cousinweb}

\begin{grammarlens}[title={What you've built}]
One more word for telling, explaining and describing at B1.
\end{grammarlens}

\subsection*{Your turn}
\begin{itemize}
\raggedright
  \item \emph{Say it:} \emph{contemplar}
  \item \emph{Say it:} \emph{contemplar}, once more
  \item \emph{Say it:} a sentence of your own with \emph{contemplar}
  \item \emph{Recall:} say the Spanish for the horizon, then the Spanish for the rainbow, then say \emph{contemplar} again
\end{itemize}

\begin{tcolorbox}[breakable,colback=teal!4,colframe=teal!35!black,title={Before you move on}]
What does contemplar mean? (\textquotedblleft{}To gaze at, to contemplate\textquotedblright{}.) Say it once more.
\end{tcolorbox}
\section[(dialogue)]{Naturaleza y tiempo}

\label{lesson:ES-R703-repaso-696-703}



\begin{quote}
Say the Spanish for the thermometer and the forecast.
\end{quote}

\subsection*{Your turn}
Say these aloud:

\begin{itemize}
\raggedright
  \item the thermometer, the forecast, the outlook, humid, the humidity
  \item to stop raining, the shower, the downpour, the drizzle, to drizzle
  \item cloudy, clear, sunny, to cloud over, to clear
  \item the lightning bolt, to flash with lightning, the hail, to hail, to thunder
  \item to snow, to freeze, the blizzard, the snowfall, the hoarfrost
  \item hot, mild, the muggy heat, the drought, to sweat
  \item to get wet, soaked, the puddle, the mud, the raincoat
  \item dusk, to get dark, the horizon, the rainbow, to gaze at
\end{itemize}

\begin{tcolorbox}[breakable,colback=teal!4,colframe=teal!35!black,title={Before you move on}]
The thermometer? (\textbf{el termómetro}.) The forecast? (\textbf{el pronóstico}.) The outlook? (\textbf{la previsión}.) Humid? (\textbf{húmedo}.) The humidity? (\textbf{la humedad}.) To stop raining? (\textbf{escampar}.) The shower? (\textbf{el chubasco}.) The downpour? (\textbf{el aguacero}.) The drizzle? (\textbf{la llovizna}.) To drizzle? (\textbf{lloviznar}.) Cloudy? (\textbf{nublado}.) Clear? (\textbf{despejado}.) Sunny? (\textbf{soleado}.) To cloud over? (\textbf{nublarse}.) To clear? (\textbf{despejar}.) The lightning bolt? (\textbf{el rayo}.) To flash with lightning? (\textbf{relampaguear}.) The hail? (\textbf{el granizo}.) To hail? (\textbf{granizar}.) To thunder? (\textbf{tronar}.) To snow? (\textbf{nevar}.) To freeze? (\textbf{helar}.) The blizzard? (\textbf{la ventisca}.) The snowfall? (\textbf{la nevada}.) The hoarfrost? (\textbf{la escarcha}.) Hot? (\textbf{caluroso}.) Mild? (\textbf{templado}.) The muggy heat? (\textbf{el bochorno}.) The drought? (\textbf{la sequía}.) To sweat? (\textbf{sudar}.) To get wet? (\textbf{mojarse}.) Soaked? (\textbf{empapado}.) The puddle? (\textbf{el charco}.) The mud? (\textbf{el barro}.) The raincoat? (\textbf{el impermeable}.) Dusk? (\textbf{el atardecer}.) To get dark? (\textbf{anochecer}.) The horizon? (\textbf{el horizonte}.) The rainbow? (\textbf{el arco iris}.) To gaze at? (\textbf{contemplar}.)
\end{tcolorbox}
\section[(dialogue)]{Al caer la tarde — once more}

\label{lesson:ES-R703-repaso-final-703}



\begin{quote}
Say the Spanish for dusk and to get dark.
\end{quote}

\subsection*{Your turn}
\begin{itemize}
\raggedright
  \item \emph{Say it:} dusk, to get dark, the horizon, the rainbow, to gaze at
\end{itemize}

\begin{tcolorbox}[breakable,colback=teal!4,colframe=teal!35!black,title={Before you move on}]
Dusk? (\textbf{el atardecer}.) To get dark? (\textbf{anochecer}.) The horizon? (\textbf{el horizonte}.) The rainbow? (\textbf{el arco iris}.) To gaze at? (\textbf{contemplar}.)
\end{tcolorbox}
\section[(dialogue)]{Lost at the station — back to the first travel words}

\label{lesson:ES-R703-repaso-problemas-de-viaje}



\begin{quote}
From the first travel chapter: say the Spanish for \textquotedblleft{}a problem\textquotedblright{} and for \textquotedblleft{}where is it?\textquotedblright{}.
\end{quote}

\subsection*{Your turn}
Say these aloud:

\begin{itemize}
\raggedright
  \item \emph{Perdone, tengo un problema con mi billete.} — and keep \emph{un}, not \emph{una}.
  \item \emph{¿Dónde está el andén cuatro?} — Where is platform four?
  \item \emph{¿Dónde está la taquilla? Necesito un billete para Madrid.}
  \item two more \emph{-ma} nouns with \emph{el}: \emph{el sistema}, \emph{el idioma}.
\end{itemize}

\begin{tcolorbox}[breakable,colback=teal!4,colframe=teal!35!black,title={Before you move on}]
Why is it \emph{el problema}? (Greek \emph{-ma} nouns came in masculine.) How do you ask where the ticket office is? (\textbf{¿Dónde está la taquilla?}) How do you ask for a ticket to Madrid? (\textbf{Un billete para Madrid, por favor}.)
\end{tcolorbox}
